% ============================================================================
% (c) 2025 Mohamed EL AFRIT -- IPSSI
% Licence CC BY-NC-SA 4.0
% https://creativecommons.org/licenses/by-nc-sa/4.0/deed.fr
%
% FICHE D'EXERCICES (enonces) -- Jour 3 : Probabilites et classification
% ============================================================================
% ============================================================================
% © 2025 Mohamed EL AFRIT — IPSSI
% Ce contenu est distribué sous licence Creative Commons BY-NC-SA 4.0
% https://creativecommons.org/licenses/by-nc-sa/4.0/deed.fr
% ============================================================================
%
% TEMPLATE COMMUN — Mathématiques appliquées à l'Intelligence Artificielle
% Ce fichier est le préambule partagé de TOUS les documents LaTeX du projet.
%
% Utilisation :
%   % ============================================================================
% © 2025 Mohamed EL AFRIT — IPSSI
% Ce contenu est distribué sous licence Creative Commons BY-NC-SA 4.0
% https://creativecommons.org/licenses/by-nc-sa/4.0/deed.fr
% ============================================================================
%
% TEMPLATE COMMUN — Mathématiques appliquées à l'Intelligence Artificielle
% Ce fichier est le préambule partagé de TOUS les documents LaTeX du projet.
%
% Utilisation :
%   % ============================================================================
% © 2025 Mohamed EL AFRIT — IPSSI
% Ce contenu est distribué sous licence Creative Commons BY-NC-SA 4.0
% https://creativecommons.org/licenses/by-nc-sa/4.0/deed.fr
% ============================================================================
%
% TEMPLATE COMMUN — Mathématiques appliquées à l'Intelligence Artificielle
% Ce fichier est le préambule partagé de TOUS les documents LaTeX du projet.
%
% Utilisation :
%   \input{template_ipssi}          % Importer ce préambule
%   \jourNumero{1}                  % Définir le numéro du jour (1, 2, 3 ou 4)
%   \jourTheme{Données et représentation mathématique}
%   \begin{document}
%   \pagegarde{Fiche de synthèse}{1}{Données et représentation mathématique}
%   ...
%   \end{document}
% ============================================================================

% --- Classe de document ---
\documentclass[11pt,a4paper]{article}

% ============================================================================
% ENCODAGE ET LANGUE
% ============================================================================
\usepackage[utf8]{inputenc}
\usepackage[T1]{fontenc}
\usepackage[french]{babel}

% ============================================================================
% MATHÉMATIQUES
% ============================================================================
\usepackage{amsmath,amssymb,amsfonts,mathtools}

% Commandes mathématiques personnalisées (conventions du cours)
\newcommand{\vx}{\mathbf{x}}          % vecteur x
\newcommand{\vw}{\mathbf{w}}          % vecteur w (poids)
\newcommand{\vy}{\mathbf{y}}          % vecteur y
\newcommand{\mX}{\mathbf{X}}          % matrice X
\newcommand{\mW}{\mathbf{W}}          % matrice W
\newcommand{\yhat}{\hat{y}}           % prédiction
\newcommand{\pscal}[2]{\langle #1, #2 \rangle} % produit scalaire
\newcommand{\gradJ}{\nabla J(\boldsymbol{\theta})} % gradient de J

% ============================================================================
% MISE EN PAGE
% ============================================================================
\usepackage[margin=2.5cm]{geometry}
\usepackage{setspace}
\setstretch{1.15}                      % interligne légèrement aéré

% ============================================================================
% COULEURS
% ============================================================================
\usepackage[dvipsnames,svgnames,x11names]{xcolor}

% Palette principale du projet
\definecolor{ipssiBleu}{HTML}{3B82F6}       % Définitions
\definecolor{ipssiVert}{HTML}{22C55E}       % À retenir
\definecolor{ipssiOrange}{HTML}{F97316}     % Attention / Piège
\definecolor{ipssiViolet}{HTML}{A855F7}     % Intuition
\definecolor{ipssiCyan}{HTML}{06B6D4}       % Exemples
\definecolor{ipssiGrisClair}{HTML}{F3F4F6}  % Fond code
\definecolor{ipssiGrisFonce}{HTML}{1F2937}  % Texte code
\definecolor{ipssiFondPage}{HTML}{FAFAFA}   % Fond léger page de garde
\definecolor{ipssiRouge}{HTML}{E74C3C}      % Classe 0 (salaire ≤ 45k€)
\definecolor{ipssiVertData}{HTML}{2ECC71}   % Classe 1 (salaire > 45k€)
\definecolor{ipssiBleuDroite}{HTML}{3498DB} % Droite de régression

% ============================================================================
% ENCADRÉS PÉDAGOGIQUES (tcolorbox)
% ============================================================================
\usepackage[most]{tcolorbox}

% Style de base partagé
\tcbset{
    ipssibase/.style={
        enhanced,
        breakable,
        left=4mm,
        right=4mm,
        top=3mm,
        bottom=3mm,
        boxrule=0pt,
        frame hidden,
        borderline west={3pt}{0pt}{#1},
        colback=#1!5,
        colframe=#1,
        fonttitle=\bfseries\sffamily,
        coltitle=#1!80!black,
        attach boxed title to top left={yshift=-2mm, xshift=4mm},
        boxed title style={
            colback=#1!15,
            colframe=#1,
            boxrule=0.5pt,
            arc=2pt,
        },
        arc=2pt,
        outer arc=2pt,
    }
}

% Icônes TikZ pour les encadrés (compatible pdflatex, pas d'emoji Unicode)
\newcommand{\icoDefinition}{%
    \tikz[baseline=-0.3em]{\draw[ipssiBleu, thick] (0,0) -- (0.3,0) -- (0.3,0.35) -- (0,0.35) -- cycle;
    \draw[ipssiBleu, thick] (0.05,0.1) -- (0.25,0.1); \draw[ipssiBleu, thick] (0.05,0.2) -- (0.2,0.2);}%
}
\newcommand{\icoRetenir}{%
    \tikz[baseline=-0.2em]{\draw[ipssiVert, thick, fill=ipssiVert!20] (0.15,0) -- (0.3,0.3) -- (0,0.3) -- cycle;}%
}
\newcommand{\icoAttention}{%
    \tikz[baseline=-0.2em]{\draw[ipssiOrange, thick] (0.15,0.35) -- (0,0) -- (0.3,0) -- cycle;
    \node[ipssiOrange, font=\tiny\bfseries] at (0.15,0.12) {!};}%
}
\newcommand{\icoIntuition}{%
    \tikz[baseline=-0.2em]{\draw[ipssiViolet, thick, fill=ipssiViolet!15]
    (0.15,0.35) -- (0.08,0.2) -- (0.08,0.15) -- (0.22,0.15) -- (0.22,0.2) -- cycle;
    \draw[ipssiViolet, thick] (0.11,0.12) -- (0.11,0.05); \draw[ipssiViolet, thick] (0.19,0.12) -- (0.19,0.05);}%
}
\newcommand{\icoExemple}{%
    \tikz[baseline=-0.2em]{\draw[ipssiCyan, thick] (0,0) rectangle (0.3,0.3);
    \draw[ipssiCyan, thick] (0,0.15) -- (0.3,0.15); \draw[ipssiCyan, thick] (0.15,0) -- (0.15,0.3);}%
}
\newcommand{\icoPython}{%
    \tikz[baseline=-0.2em]{\node[font=\small\bfseries\ttfamily, ipssiVert!70!black] at (0,0) {\textgreater\textgreater};}%
}

% Définition (bleu)
\newtcolorbox{definitionbox}[1][D\'{e}finition]{
    ipssibase=ipssiBleu,
    title={\icoDefinition\; #1},
}

% À retenir (vert)
\newtcolorbox{retenirbox}[1][\`{A} retenir]{
    ipssibase=ipssiVert,
    title={\icoRetenir\; #1},
}

% Attention / Piège (orange)
\newtcolorbox{attentionbox}[1][Attention]{
    ipssibase=ipssiOrange,
    title={\icoAttention\; #1},
}

% Intuition (violet)
\newtcolorbox{intuitionbox}[1][Intuition]{
    ipssibase=ipssiViolet,
    title={\icoIntuition\; #1},
}

% Exemple (cyan)
\newtcolorbox{exemplebox}[1][Exemple]{
    ipssibase=ipssiCyan,
    title={\icoExemple\; #1},
}

% Code Python (style terminal)
\newtcolorbox{pythonbox}[1][Code Python]{
    enhanced,
    breakable,
    left=4mm, right=4mm, top=3mm, bottom=3mm,
    boxrule=0.5pt,
    colback=ipssiGrisClair,
    colframe=ipssiGrisFonce!40,
    fonttitle=\bfseries\ttfamily\small,
    coltitle=ipssiGrisFonce,
    title={\icoPython\; #1},
    attach boxed title to top left={yshift=-2mm, xshift=4mm},
    boxed title style={colback=ipssiGrisClair, colframe=ipssiGrisFonce!40, boxrule=0.5pt, arc=2pt},
    arc=2pt, outer arc=2pt,
}

% ============================================================================
% CODE PYTHON (listings)
% ============================================================================
\usepackage{listings}

\lstdefinestyle{python}{
    language=Python,
    basicstyle=\ttfamily\small,
    keywordstyle=\color{ipssiBleu}\bfseries,
    stringstyle=\color{ipssiVert!70!black},
    commentstyle=\color{gray}\itshape,
    numberstyle=\tiny\color{gray},
    numbers=left,
    numbersep=8pt,
    backgroundcolor=\color{ipssiGrisClair},
    frame=none,
    rulecolor=\color{ipssiGrisFonce!20},
    breaklines=true,
    breakatwhitespace=true,
    tabsize=4,
    showstringspaces=false,
    captionpos=b,
    aboveskip=8pt,
    belowskip=8pt,
    xleftmargin=10pt,
    framexleftmargin=10pt,
    morekeywords={as, with, True, False, None, self, np, plt, import, from},
    literate=
        {é}{{\'e}}1 {è}{{\`e}}1 {ê}{{\^e}}1 {ë}{{\"e}}1
        {à}{{\`a}}1 {â}{{\^a}}1
        {ù}{{\`u}}1 {û}{{\^u}}1
        {ô}{{\^o}}1
        {î}{{\^i}}1 {ï}{{\"i}}1
        {ç}{{\c{c}}}1
        {É}{{\'E}}1 {È}{{\`E}}1
        {→}{{$\rightarrow$}}1
        {≈}{{$\approx$}}1
        {≤}{{$\leq$}}1
        {≥}{{$\geq$}}1
        {★}{{$\bigstar$}}1
        {ŷ}{{$\hat{y}$}}1
        {·}{{$\cdot$}}1,
}

\lstset{style=python}

% Style pour le code inline
\newcommand{\pyinline}[1]{\lstinline[style=python,basicstyle=\ttfamily\small]{#1}}

% ============================================================================
% GRAPHIQUES
% ============================================================================
\usepackage{tikz}
\usepackage{pgfplots}
\pgfplotsset{compat=1.18}
\usetikzlibrary{arrows.meta, calc, positioning, shapes.geometric, decorations.pathreplacing}

% ============================================================================
% TABLEAUX
% ============================================================================
\usepackage{booktabs}
\usepackage{tabularx}
\usepackage{multirow}

% ============================================================================
% LISTES
% ============================================================================
\usepackage{enumitem}
\setlist[itemize]{itemsep=2pt, parsep=0pt, topsep=4pt}
\setlist[enumerate]{itemsep=2pt, parsep=0pt, topsep=4pt}

% ============================================================================
% HYPERLIENS
% ============================================================================
\usepackage[
    colorlinks=true,
    linkcolor=ipssiBleu!80!black,
    urlcolor=ipssiBleu!80!black,
    citecolor=ipssiVert!80!black,
    bookmarks=true,
    bookmarksopen=true,
    pdfauthor={Mohamed EL AFRIT},
    pdftitle={Mathématiques appliquées à l'IA — IPSSI},
    pdfsubject={Formation Bachelor 2 — IPSSI 2025-2026},
]{hyperref}

% ============================================================================
% IMAGES
% ============================================================================
\usepackage{graphicx}

% ============================================================================
% VARIABLES PARAMÉTRABLES (jour et thème)
% ============================================================================
\newcommand{\leJour}{X}
\newcommand{\leTheme}{Thème}
\newcommand{\jourNumero}[1]{\renewcommand{\leJour}{#1}}
\newcommand{\jourTheme}[1]{\renewcommand{\leTheme}{#1}}

% ============================================================================
% EN-TÊTE ET PIED DE PAGE (fancyhdr)
% ============================================================================
\usepackage{fancyhdr}

\pagestyle{fancy}
\fancyhf{}

% En-tête
\fancyhead[L]{\small\sffamily\textcolor{gray}{IPSSI — Maths appliquées à l'IA}}
\fancyhead[R]{\small\sffamily\textcolor{gray}{Jour \leJour\ — \leTheme}}
\renewcommand{\headrulewidth}{0.4pt}

% Pied de page
\fancyfoot[L]{\footnotesize\textcolor{gray}{Mohamed EL AFRIT\enspace\textbullet\enspace\url{www.mohamedelafrit.com}\enspace\textbullet\enspace CC BY-NC-SA 4.0}}
\fancyfoot[C]{\footnotesize\textcolor{gray}{\thepage}}
\fancyfoot[R]{\footnotesize\textcolor{gray}{2025–2026}}
\renewcommand{\footrulewidth}{0.2pt}

% Style pour la page de garde (pas d'en-tête)
\fancypagestyle{pagegarde}{
    \fancyhf{}
    \renewcommand{\headrulewidth}{0pt}
    \fancyfoot[C]{\footnotesize\textcolor{gray}{\thepage}}
    \renewcommand{\footrulewidth}{0pt}
}

% ============================================================================
% PAGE DE GARDE — \pagegarde{titre}{jour}{theme}
% ============================================================================
\newcommand{\pagegarde}[3]{%
    \thispagestyle{pagegarde}
    \begin{center}
        \vspace*{1.5cm}

        % Logo / Bandeau institution
        {\Large\sffamily\bfseries\textcolor{ipssiBleu}{IPSSI}}\par
        \vspace{2pt}
        {\normalsize\sffamily 2\textsuperscript{e} année Bachelor Informatique}\par

        \vspace{0.8cm}
        {\color{ipssiBleu}\rule{0.6\textwidth}{1.5pt}}\par
        \vspace{0.8cm}

        % Titre du module
        {\LARGE\sffamily\bfseries Mathématiques appliquées\\[4pt] à l'Intelligence Artificielle}\par

        \vspace{1cm}

        % Titre du document
        \begin{tcolorbox}[
            enhanced,
            colback=ipssiBleu!5,
            colframe=ipssiBleu!50,
            boxrule=1pt,
            arc=4pt,
            width=0.85\textwidth,
            halign=center,
            top=10pt, bottom=10pt,
        ]
            {\huge\sffamily\bfseries\textcolor{ipssiBleu!80!black}{#1}}\par
            \vspace{6pt}
            {\Large\sffamily Jour #2 — #3}
        \end{tcolorbox}

        \vspace{1.5cm}

        % Informations auteur
        {\large\sffamily
            \textbf{Auteur :} Mohamed EL AFRIT\\[4pt]
            \url{www.mohamedelafrit.com}\\[8pt]
            \textbf{Année académique :} 2025–2026
        }\par

        \vspace{1.5cm}

        % Licence Creative Commons
        {\color{ipssiBleu}\rule{0.4\textwidth}{0.5pt}}\par
        \vspace{8pt}
        {\small\sffamily
            \textcopyright\ 2025 Mohamed EL AFRIT — IPSSI\\[2pt]
            Ce contenu est distribué sous licence\\[2pt]
            \textbf{Creative Commons BY-NC-SA 4.0}\\[2pt]
            \url{https://creativecommons.org/licenses/by-nc-sa/4.0/deed.fr}
        }\par
    \end{center}
    \newpage
}

% ============================================================================
% NIVEAUX D'EXERCICES
% ============================================================================

% Étoile compatible pdflatex
\newcommand{\starfull}[1]{\textcolor{#1}{$\bigstar$}}
\newcommand{\starempty}{\textcolor{gray!40}{$\bigstar$}}

% Fondamental (1 étoile)
\newcommand{\exofondamental}{%
    \starfull{ipssiVert}\starempty\starempty\enspace%
    {\small\sffamily\textcolor{ipssiVert}{\textbf{Fondamental}}}%
}

% Intermédiaire (2 étoiles)
\newcommand{\exointermediaire}{%
    \starfull{ipssiOrange}\starfull{ipssiOrange}\starempty\enspace%
    {\small\sffamily\textcolor{ipssiOrange}{\textbf{Interm\'{e}diaire}}}%
}

% Avancé (3 étoiles)
\newcommand{\exoavance}{%
    \starfull{ipssiRouge}\starfull{ipssiRouge}\starfull{ipssiRouge}\enspace%
    {\small\sffamily\textcolor{ipssiRouge}{\textbf{Avanc\'{e}}}}%
}

% Commande d'exercice numéroté
\newcounter{exocount}[section]
\newcommand{\exercice}[1]{%
    \stepcounter{exocount}%
    \vspace{8pt}%
    \noindent\textbf{\sffamily Exercice \theexocount}\enspace — \enspace #1\par
    \vspace{4pt}%
}

% ============================================================================
% COMMANDES UTILITAIRES
% ============================================================================

% Séparateur visuel entre sections
\newcommand{\separateur}{%
    \vspace{6pt}
    \begin{center}
        {\color{ipssiBleu!30}\rule{0.3\textwidth}{0.5pt}}
    \end{center}
    \vspace{6pt}
}

% Mot-clé en gras coloré
\newcommand{\motcle}[1]{\textbf{\textcolor{ipssiBleu!80!black}{#1}}}

% Formule encadrée importante
\newcommand{\formule}[1]{%
    \begin{center}
        \fcolorbox{ipssiBleu!50}{ipssiBleu!5}{%
            \quad$\displaystyle #1$\quad%
        }
    \end{center}
}

% ============================================================================
% FIN DU PRÉAMBULE
% ============================================================================
          % Importer ce préambule
%   \jourNumero{1}                  % Définir le numéro du jour (1, 2, 3 ou 4)
%   \jourTheme{Données et représentation mathématique}
%   \begin{document}
%   \pagegarde{Fiche de synthèse}{1}{Données et représentation mathématique}
%   ...
%   \end{document}
% ============================================================================

% --- Classe de document ---
\documentclass[11pt,a4paper]{article}

% ============================================================================
% ENCODAGE ET LANGUE
% ============================================================================
\usepackage[utf8]{inputenc}
\usepackage[T1]{fontenc}
\usepackage[french]{babel}

% ============================================================================
% MATHÉMATIQUES
% ============================================================================
\usepackage{amsmath,amssymb,amsfonts,mathtools}

% Commandes mathématiques personnalisées (conventions du cours)
\newcommand{\vx}{\mathbf{x}}          % vecteur x
\newcommand{\vw}{\mathbf{w}}          % vecteur w (poids)
\newcommand{\vy}{\mathbf{y}}          % vecteur y
\newcommand{\mX}{\mathbf{X}}          % matrice X
\newcommand{\mW}{\mathbf{W}}          % matrice W
\newcommand{\yhat}{\hat{y}}           % prédiction
\newcommand{\pscal}[2]{\langle #1, #2 \rangle} % produit scalaire
\newcommand{\gradJ}{\nabla J(\boldsymbol{\theta})} % gradient de J

% ============================================================================
% MISE EN PAGE
% ============================================================================
\usepackage[margin=2.5cm]{geometry}
\usepackage{setspace}
\setstretch{1.15}                      % interligne légèrement aéré

% ============================================================================
% COULEURS
% ============================================================================
\usepackage[dvipsnames,svgnames,x11names]{xcolor}

% Palette principale du projet
\definecolor{ipssiBleu}{HTML}{3B82F6}       % Définitions
\definecolor{ipssiVert}{HTML}{22C55E}       % À retenir
\definecolor{ipssiOrange}{HTML}{F97316}     % Attention / Piège
\definecolor{ipssiViolet}{HTML}{A855F7}     % Intuition
\definecolor{ipssiCyan}{HTML}{06B6D4}       % Exemples
\definecolor{ipssiGrisClair}{HTML}{F3F4F6}  % Fond code
\definecolor{ipssiGrisFonce}{HTML}{1F2937}  % Texte code
\definecolor{ipssiFondPage}{HTML}{FAFAFA}   % Fond léger page de garde
\definecolor{ipssiRouge}{HTML}{E74C3C}      % Classe 0 (salaire ≤ 45k€)
\definecolor{ipssiVertData}{HTML}{2ECC71}   % Classe 1 (salaire > 45k€)
\definecolor{ipssiBleuDroite}{HTML}{3498DB} % Droite de régression

% ============================================================================
% ENCADRÉS PÉDAGOGIQUES (tcolorbox)
% ============================================================================
\usepackage[most]{tcolorbox}

% Style de base partagé
\tcbset{
    ipssibase/.style={
        enhanced,
        breakable,
        left=4mm,
        right=4mm,
        top=3mm,
        bottom=3mm,
        boxrule=0pt,
        frame hidden,
        borderline west={3pt}{0pt}{#1},
        colback=#1!5,
        colframe=#1,
        fonttitle=\bfseries\sffamily,
        coltitle=#1!80!black,
        attach boxed title to top left={yshift=-2mm, xshift=4mm},
        boxed title style={
            colback=#1!15,
            colframe=#1,
            boxrule=0.5pt,
            arc=2pt,
        },
        arc=2pt,
        outer arc=2pt,
    }
}

% Icônes TikZ pour les encadrés (compatible pdflatex, pas d'emoji Unicode)
\newcommand{\icoDefinition}{%
    \tikz[baseline=-0.3em]{\draw[ipssiBleu, thick] (0,0) -- (0.3,0) -- (0.3,0.35) -- (0,0.35) -- cycle;
    \draw[ipssiBleu, thick] (0.05,0.1) -- (0.25,0.1); \draw[ipssiBleu, thick] (0.05,0.2) -- (0.2,0.2);}%
}
\newcommand{\icoRetenir}{%
    \tikz[baseline=-0.2em]{\draw[ipssiVert, thick, fill=ipssiVert!20] (0.15,0) -- (0.3,0.3) -- (0,0.3) -- cycle;}%
}
\newcommand{\icoAttention}{%
    \tikz[baseline=-0.2em]{\draw[ipssiOrange, thick] (0.15,0.35) -- (0,0) -- (0.3,0) -- cycle;
    \node[ipssiOrange, font=\tiny\bfseries] at (0.15,0.12) {!};}%
}
\newcommand{\icoIntuition}{%
    \tikz[baseline=-0.2em]{\draw[ipssiViolet, thick, fill=ipssiViolet!15]
    (0.15,0.35) -- (0.08,0.2) -- (0.08,0.15) -- (0.22,0.15) -- (0.22,0.2) -- cycle;
    \draw[ipssiViolet, thick] (0.11,0.12) -- (0.11,0.05); \draw[ipssiViolet, thick] (0.19,0.12) -- (0.19,0.05);}%
}
\newcommand{\icoExemple}{%
    \tikz[baseline=-0.2em]{\draw[ipssiCyan, thick] (0,0) rectangle (0.3,0.3);
    \draw[ipssiCyan, thick] (0,0.15) -- (0.3,0.15); \draw[ipssiCyan, thick] (0.15,0) -- (0.15,0.3);}%
}
\newcommand{\icoPython}{%
    \tikz[baseline=-0.2em]{\node[font=\small\bfseries\ttfamily, ipssiVert!70!black] at (0,0) {\textgreater\textgreater};}%
}

% Définition (bleu)
\newtcolorbox{definitionbox}[1][D\'{e}finition]{
    ipssibase=ipssiBleu,
    title={\icoDefinition\; #1},
}

% À retenir (vert)
\newtcolorbox{retenirbox}[1][\`{A} retenir]{
    ipssibase=ipssiVert,
    title={\icoRetenir\; #1},
}

% Attention / Piège (orange)
\newtcolorbox{attentionbox}[1][Attention]{
    ipssibase=ipssiOrange,
    title={\icoAttention\; #1},
}

% Intuition (violet)
\newtcolorbox{intuitionbox}[1][Intuition]{
    ipssibase=ipssiViolet,
    title={\icoIntuition\; #1},
}

% Exemple (cyan)
\newtcolorbox{exemplebox}[1][Exemple]{
    ipssibase=ipssiCyan,
    title={\icoExemple\; #1},
}

% Code Python (style terminal)
\newtcolorbox{pythonbox}[1][Code Python]{
    enhanced,
    breakable,
    left=4mm, right=4mm, top=3mm, bottom=3mm,
    boxrule=0.5pt,
    colback=ipssiGrisClair,
    colframe=ipssiGrisFonce!40,
    fonttitle=\bfseries\ttfamily\small,
    coltitle=ipssiGrisFonce,
    title={\icoPython\; #1},
    attach boxed title to top left={yshift=-2mm, xshift=4mm},
    boxed title style={colback=ipssiGrisClair, colframe=ipssiGrisFonce!40, boxrule=0.5pt, arc=2pt},
    arc=2pt, outer arc=2pt,
}

% ============================================================================
% CODE PYTHON (listings)
% ============================================================================
\usepackage{listings}

\lstdefinestyle{python}{
    language=Python,
    basicstyle=\ttfamily\small,
    keywordstyle=\color{ipssiBleu}\bfseries,
    stringstyle=\color{ipssiVert!70!black},
    commentstyle=\color{gray}\itshape,
    numberstyle=\tiny\color{gray},
    numbers=left,
    numbersep=8pt,
    backgroundcolor=\color{ipssiGrisClair},
    frame=none,
    rulecolor=\color{ipssiGrisFonce!20},
    breaklines=true,
    breakatwhitespace=true,
    tabsize=4,
    showstringspaces=false,
    captionpos=b,
    aboveskip=8pt,
    belowskip=8pt,
    xleftmargin=10pt,
    framexleftmargin=10pt,
    morekeywords={as, with, True, False, None, self, np, plt, import, from},
    literate=
        {é}{{\'e}}1 {è}{{\`e}}1 {ê}{{\^e}}1 {ë}{{\"e}}1
        {à}{{\`a}}1 {â}{{\^a}}1
        {ù}{{\`u}}1 {û}{{\^u}}1
        {ô}{{\^o}}1
        {î}{{\^i}}1 {ï}{{\"i}}1
        {ç}{{\c{c}}}1
        {É}{{\'E}}1 {È}{{\`E}}1
        {→}{{$\rightarrow$}}1
        {≈}{{$\approx$}}1
        {≤}{{$\leq$}}1
        {≥}{{$\geq$}}1
        {★}{{$\bigstar$}}1
        {ŷ}{{$\hat{y}$}}1
        {·}{{$\cdot$}}1,
}

\lstset{style=python}

% Style pour le code inline
\newcommand{\pyinline}[1]{\lstinline[style=python,basicstyle=\ttfamily\small]{#1}}

% ============================================================================
% GRAPHIQUES
% ============================================================================
\usepackage{tikz}
\usepackage{pgfplots}
\pgfplotsset{compat=1.18}
\usetikzlibrary{arrows.meta, calc, positioning, shapes.geometric, decorations.pathreplacing}

% ============================================================================
% TABLEAUX
% ============================================================================
\usepackage{booktabs}
\usepackage{tabularx}
\usepackage{multirow}

% ============================================================================
% LISTES
% ============================================================================
\usepackage{enumitem}
\setlist[itemize]{itemsep=2pt, parsep=0pt, topsep=4pt}
\setlist[enumerate]{itemsep=2pt, parsep=0pt, topsep=4pt}

% ============================================================================
% HYPERLIENS
% ============================================================================
\usepackage[
    colorlinks=true,
    linkcolor=ipssiBleu!80!black,
    urlcolor=ipssiBleu!80!black,
    citecolor=ipssiVert!80!black,
    bookmarks=true,
    bookmarksopen=true,
    pdfauthor={Mohamed EL AFRIT},
    pdftitle={Mathématiques appliquées à l'IA — IPSSI},
    pdfsubject={Formation Bachelor 2 — IPSSI 2025-2026},
]{hyperref}

% ============================================================================
% IMAGES
% ============================================================================
\usepackage{graphicx}

% ============================================================================
% VARIABLES PARAMÉTRABLES (jour et thème)
% ============================================================================
\newcommand{\leJour}{X}
\newcommand{\leTheme}{Thème}
\newcommand{\jourNumero}[1]{\renewcommand{\leJour}{#1}}
\newcommand{\jourTheme}[1]{\renewcommand{\leTheme}{#1}}

% ============================================================================
% EN-TÊTE ET PIED DE PAGE (fancyhdr)
% ============================================================================
\usepackage{fancyhdr}

\pagestyle{fancy}
\fancyhf{}

% En-tête
\fancyhead[L]{\small\sffamily\textcolor{gray}{IPSSI — Maths appliquées à l'IA}}
\fancyhead[R]{\small\sffamily\textcolor{gray}{Jour \leJour\ — \leTheme}}
\renewcommand{\headrulewidth}{0.4pt}

% Pied de page
\fancyfoot[L]{\footnotesize\textcolor{gray}{Mohamed EL AFRIT\enspace\textbullet\enspace\url{www.mohamedelafrit.com}\enspace\textbullet\enspace CC BY-NC-SA 4.0}}
\fancyfoot[C]{\footnotesize\textcolor{gray}{\thepage}}
\fancyfoot[R]{\footnotesize\textcolor{gray}{2025–2026}}
\renewcommand{\footrulewidth}{0.2pt}

% Style pour la page de garde (pas d'en-tête)
\fancypagestyle{pagegarde}{
    \fancyhf{}
    \renewcommand{\headrulewidth}{0pt}
    \fancyfoot[C]{\footnotesize\textcolor{gray}{\thepage}}
    \renewcommand{\footrulewidth}{0pt}
}

% ============================================================================
% PAGE DE GARDE — \pagegarde{titre}{jour}{theme}
% ============================================================================
\newcommand{\pagegarde}[3]{%
    \thispagestyle{pagegarde}
    \begin{center}
        \vspace*{1.5cm}

        % Logo / Bandeau institution
        {\Large\sffamily\bfseries\textcolor{ipssiBleu}{IPSSI}}\par
        \vspace{2pt}
        {\normalsize\sffamily 2\textsuperscript{e} année Bachelor Informatique}\par

        \vspace{0.8cm}
        {\color{ipssiBleu}\rule{0.6\textwidth}{1.5pt}}\par
        \vspace{0.8cm}

        % Titre du module
        {\LARGE\sffamily\bfseries Mathématiques appliquées\\[4pt] à l'Intelligence Artificielle}\par

        \vspace{1cm}

        % Titre du document
        \begin{tcolorbox}[
            enhanced,
            colback=ipssiBleu!5,
            colframe=ipssiBleu!50,
            boxrule=1pt,
            arc=4pt,
            width=0.85\textwidth,
            halign=center,
            top=10pt, bottom=10pt,
        ]
            {\huge\sffamily\bfseries\textcolor{ipssiBleu!80!black}{#1}}\par
            \vspace{6pt}
            {\Large\sffamily Jour #2 — #3}
        \end{tcolorbox}

        \vspace{1.5cm}

        % Informations auteur
        {\large\sffamily
            \textbf{Auteur :} Mohamed EL AFRIT\\[4pt]
            \url{www.mohamedelafrit.com}\\[8pt]
            \textbf{Année académique :} 2025–2026
        }\par

        \vspace{1.5cm}

        % Licence Creative Commons
        {\color{ipssiBleu}\rule{0.4\textwidth}{0.5pt}}\par
        \vspace{8pt}
        {\small\sffamily
            \textcopyright\ 2025 Mohamed EL AFRIT — IPSSI\\[2pt]
            Ce contenu est distribué sous licence\\[2pt]
            \textbf{Creative Commons BY-NC-SA 4.0}\\[2pt]
            \url{https://creativecommons.org/licenses/by-nc-sa/4.0/deed.fr}
        }\par
    \end{center}
    \newpage
}

% ============================================================================
% NIVEAUX D'EXERCICES
% ============================================================================

% Étoile compatible pdflatex
\newcommand{\starfull}[1]{\textcolor{#1}{$\bigstar$}}
\newcommand{\starempty}{\textcolor{gray!40}{$\bigstar$}}

% Fondamental (1 étoile)
\newcommand{\exofondamental}{%
    \starfull{ipssiVert}\starempty\starempty\enspace%
    {\small\sffamily\textcolor{ipssiVert}{\textbf{Fondamental}}}%
}

% Intermédiaire (2 étoiles)
\newcommand{\exointermediaire}{%
    \starfull{ipssiOrange}\starfull{ipssiOrange}\starempty\enspace%
    {\small\sffamily\textcolor{ipssiOrange}{\textbf{Interm\'{e}diaire}}}%
}

% Avancé (3 étoiles)
\newcommand{\exoavance}{%
    \starfull{ipssiRouge}\starfull{ipssiRouge}\starfull{ipssiRouge}\enspace%
    {\small\sffamily\textcolor{ipssiRouge}{\textbf{Avanc\'{e}}}}%
}

% Commande d'exercice numéroté
\newcounter{exocount}[section]
\newcommand{\exercice}[1]{%
    \stepcounter{exocount}%
    \vspace{8pt}%
    \noindent\textbf{\sffamily Exercice \theexocount}\enspace — \enspace #1\par
    \vspace{4pt}%
}

% ============================================================================
% COMMANDES UTILITAIRES
% ============================================================================

% Séparateur visuel entre sections
\newcommand{\separateur}{%
    \vspace{6pt}
    \begin{center}
        {\color{ipssiBleu!30}\rule{0.3\textwidth}{0.5pt}}
    \end{center}
    \vspace{6pt}
}

% Mot-clé en gras coloré
\newcommand{\motcle}[1]{\textbf{\textcolor{ipssiBleu!80!black}{#1}}}

% Formule encadrée importante
\newcommand{\formule}[1]{%
    \begin{center}
        \fcolorbox{ipssiBleu!50}{ipssiBleu!5}{%
            \quad$\displaystyle #1$\quad%
        }
    \end{center}
}

% ============================================================================
% FIN DU PRÉAMBULE
% ============================================================================
          % Importer ce préambule
%   \jourNumero{1}                  % Définir le numéro du jour (1, 2, 3 ou 4)
%   \jourTheme{Données et représentation mathématique}
%   \begin{document}
%   \pagegarde{Fiche de synthèse}{1}{Données et représentation mathématique}
%   ...
%   \end{document}
% ============================================================================

% --- Classe de document ---
\documentclass[11pt,a4paper]{article}

% ============================================================================
% ENCODAGE ET LANGUE
% ============================================================================
\usepackage[utf8]{inputenc}
\usepackage[T1]{fontenc}
\usepackage[french]{babel}

% ============================================================================
% MATHÉMATIQUES
% ============================================================================
\usepackage{amsmath,amssymb,amsfonts,mathtools}

% Commandes mathématiques personnalisées (conventions du cours)
\newcommand{\vx}{\mathbf{x}}          % vecteur x
\newcommand{\vw}{\mathbf{w}}          % vecteur w (poids)
\newcommand{\vy}{\mathbf{y}}          % vecteur y
\newcommand{\mX}{\mathbf{X}}          % matrice X
\newcommand{\mW}{\mathbf{W}}          % matrice W
\newcommand{\yhat}{\hat{y}}           % prédiction
\newcommand{\pscal}[2]{\langle #1, #2 \rangle} % produit scalaire
\newcommand{\gradJ}{\nabla J(\boldsymbol{\theta})} % gradient de J

% ============================================================================
% MISE EN PAGE
% ============================================================================
\usepackage[margin=2.5cm]{geometry}
\usepackage{setspace}
\setstretch{1.15}                      % interligne légèrement aéré

% ============================================================================
% COULEURS
% ============================================================================
\usepackage[dvipsnames,svgnames,x11names]{xcolor}

% Palette principale du projet
\definecolor{ipssiBleu}{HTML}{3B82F6}       % Définitions
\definecolor{ipssiVert}{HTML}{22C55E}       % À retenir
\definecolor{ipssiOrange}{HTML}{F97316}     % Attention / Piège
\definecolor{ipssiViolet}{HTML}{A855F7}     % Intuition
\definecolor{ipssiCyan}{HTML}{06B6D4}       % Exemples
\definecolor{ipssiGrisClair}{HTML}{F3F4F6}  % Fond code
\definecolor{ipssiGrisFonce}{HTML}{1F2937}  % Texte code
\definecolor{ipssiFondPage}{HTML}{FAFAFA}   % Fond léger page de garde
\definecolor{ipssiRouge}{HTML}{E74C3C}      % Classe 0 (salaire ≤ 45k€)
\definecolor{ipssiVertData}{HTML}{2ECC71}   % Classe 1 (salaire > 45k€)
\definecolor{ipssiBleuDroite}{HTML}{3498DB} % Droite de régression

% ============================================================================
% ENCADRÉS PÉDAGOGIQUES (tcolorbox)
% ============================================================================
\usepackage[most]{tcolorbox}

% Style de base partagé
\tcbset{
    ipssibase/.style={
        enhanced,
        breakable,
        left=4mm,
        right=4mm,
        top=3mm,
        bottom=3mm,
        boxrule=0pt,
        frame hidden,
        borderline west={3pt}{0pt}{#1},
        colback=#1!5,
        colframe=#1,
        fonttitle=\bfseries\sffamily,
        coltitle=#1!80!black,
        attach boxed title to top left={yshift=-2mm, xshift=4mm},
        boxed title style={
            colback=#1!15,
            colframe=#1,
            boxrule=0.5pt,
            arc=2pt,
        },
        arc=2pt,
        outer arc=2pt,
    }
}

% Icônes TikZ pour les encadrés (compatible pdflatex, pas d'emoji Unicode)
\newcommand{\icoDefinition}{%
    \tikz[baseline=-0.3em]{\draw[ipssiBleu, thick] (0,0) -- (0.3,0) -- (0.3,0.35) -- (0,0.35) -- cycle;
    \draw[ipssiBleu, thick] (0.05,0.1) -- (0.25,0.1); \draw[ipssiBleu, thick] (0.05,0.2) -- (0.2,0.2);}%
}
\newcommand{\icoRetenir}{%
    \tikz[baseline=-0.2em]{\draw[ipssiVert, thick, fill=ipssiVert!20] (0.15,0) -- (0.3,0.3) -- (0,0.3) -- cycle;}%
}
\newcommand{\icoAttention}{%
    \tikz[baseline=-0.2em]{\draw[ipssiOrange, thick] (0.15,0.35) -- (0,0) -- (0.3,0) -- cycle;
    \node[ipssiOrange, font=\tiny\bfseries] at (0.15,0.12) {!};}%
}
\newcommand{\icoIntuition}{%
    \tikz[baseline=-0.2em]{\draw[ipssiViolet, thick, fill=ipssiViolet!15]
    (0.15,0.35) -- (0.08,0.2) -- (0.08,0.15) -- (0.22,0.15) -- (0.22,0.2) -- cycle;
    \draw[ipssiViolet, thick] (0.11,0.12) -- (0.11,0.05); \draw[ipssiViolet, thick] (0.19,0.12) -- (0.19,0.05);}%
}
\newcommand{\icoExemple}{%
    \tikz[baseline=-0.2em]{\draw[ipssiCyan, thick] (0,0) rectangle (0.3,0.3);
    \draw[ipssiCyan, thick] (0,0.15) -- (0.3,0.15); \draw[ipssiCyan, thick] (0.15,0) -- (0.15,0.3);}%
}
\newcommand{\icoPython}{%
    \tikz[baseline=-0.2em]{\node[font=\small\bfseries\ttfamily, ipssiVert!70!black] at (0,0) {\textgreater\textgreater};}%
}

% Définition (bleu)
\newtcolorbox{definitionbox}[1][D\'{e}finition]{
    ipssibase=ipssiBleu,
    title={\icoDefinition\; #1},
}

% À retenir (vert)
\newtcolorbox{retenirbox}[1][\`{A} retenir]{
    ipssibase=ipssiVert,
    title={\icoRetenir\; #1},
}

% Attention / Piège (orange)
\newtcolorbox{attentionbox}[1][Attention]{
    ipssibase=ipssiOrange,
    title={\icoAttention\; #1},
}

% Intuition (violet)
\newtcolorbox{intuitionbox}[1][Intuition]{
    ipssibase=ipssiViolet,
    title={\icoIntuition\; #1},
}

% Exemple (cyan)
\newtcolorbox{exemplebox}[1][Exemple]{
    ipssibase=ipssiCyan,
    title={\icoExemple\; #1},
}

% Code Python (style terminal)
\newtcolorbox{pythonbox}[1][Code Python]{
    enhanced,
    breakable,
    left=4mm, right=4mm, top=3mm, bottom=3mm,
    boxrule=0.5pt,
    colback=ipssiGrisClair,
    colframe=ipssiGrisFonce!40,
    fonttitle=\bfseries\ttfamily\small,
    coltitle=ipssiGrisFonce,
    title={\icoPython\; #1},
    attach boxed title to top left={yshift=-2mm, xshift=4mm},
    boxed title style={colback=ipssiGrisClair, colframe=ipssiGrisFonce!40, boxrule=0.5pt, arc=2pt},
    arc=2pt, outer arc=2pt,
}

% ============================================================================
% CODE PYTHON (listings)
% ============================================================================
\usepackage{listings}

\lstdefinestyle{python}{
    language=Python,
    basicstyle=\ttfamily\small,
    keywordstyle=\color{ipssiBleu}\bfseries,
    stringstyle=\color{ipssiVert!70!black},
    commentstyle=\color{gray}\itshape,
    numberstyle=\tiny\color{gray},
    numbers=left,
    numbersep=8pt,
    backgroundcolor=\color{ipssiGrisClair},
    frame=none,
    rulecolor=\color{ipssiGrisFonce!20},
    breaklines=true,
    breakatwhitespace=true,
    tabsize=4,
    showstringspaces=false,
    captionpos=b,
    aboveskip=8pt,
    belowskip=8pt,
    xleftmargin=10pt,
    framexleftmargin=10pt,
    morekeywords={as, with, True, False, None, self, np, plt, import, from},
    literate=
        {é}{{\'e}}1 {è}{{\`e}}1 {ê}{{\^e}}1 {ë}{{\"e}}1
        {à}{{\`a}}1 {â}{{\^a}}1
        {ù}{{\`u}}1 {û}{{\^u}}1
        {ô}{{\^o}}1
        {î}{{\^i}}1 {ï}{{\"i}}1
        {ç}{{\c{c}}}1
        {É}{{\'E}}1 {È}{{\`E}}1
        {→}{{$\rightarrow$}}1
        {≈}{{$\approx$}}1
        {≤}{{$\leq$}}1
        {≥}{{$\geq$}}1
        {★}{{$\bigstar$}}1
        {ŷ}{{$\hat{y}$}}1
        {·}{{$\cdot$}}1,
}

\lstset{style=python}

% Style pour le code inline
\newcommand{\pyinline}[1]{\lstinline[style=python,basicstyle=\ttfamily\small]{#1}}

% ============================================================================
% GRAPHIQUES
% ============================================================================
\usepackage{tikz}
\usepackage{pgfplots}
\pgfplotsset{compat=1.18}
\usetikzlibrary{arrows.meta, calc, positioning, shapes.geometric, decorations.pathreplacing}

% ============================================================================
% TABLEAUX
% ============================================================================
\usepackage{booktabs}
\usepackage{tabularx}
\usepackage{multirow}

% ============================================================================
% LISTES
% ============================================================================
\usepackage{enumitem}
\setlist[itemize]{itemsep=2pt, parsep=0pt, topsep=4pt}
\setlist[enumerate]{itemsep=2pt, parsep=0pt, topsep=4pt}

% ============================================================================
% HYPERLIENS
% ============================================================================
\usepackage[
    colorlinks=true,
    linkcolor=ipssiBleu!80!black,
    urlcolor=ipssiBleu!80!black,
    citecolor=ipssiVert!80!black,
    bookmarks=true,
    bookmarksopen=true,
    pdfauthor={Mohamed EL AFRIT},
    pdftitle={Mathématiques appliquées à l'IA — IPSSI},
    pdfsubject={Formation Bachelor 2 — IPSSI 2025-2026},
]{hyperref}

% ============================================================================
% IMAGES
% ============================================================================
\usepackage{graphicx}

% ============================================================================
% VARIABLES PARAMÉTRABLES (jour et thème)
% ============================================================================
\newcommand{\leJour}{X}
\newcommand{\leTheme}{Thème}
\newcommand{\jourNumero}[1]{\renewcommand{\leJour}{#1}}
\newcommand{\jourTheme}[1]{\renewcommand{\leTheme}{#1}}

% ============================================================================
% EN-TÊTE ET PIED DE PAGE (fancyhdr)
% ============================================================================
\usepackage{fancyhdr}

\pagestyle{fancy}
\fancyhf{}

% En-tête
\fancyhead[L]{\small\sffamily\textcolor{gray}{IPSSI — Maths appliquées à l'IA}}
\fancyhead[R]{\small\sffamily\textcolor{gray}{Jour \leJour\ — \leTheme}}
\renewcommand{\headrulewidth}{0.4pt}

% Pied de page
\fancyfoot[L]{\footnotesize\textcolor{gray}{Mohamed EL AFRIT\enspace\textbullet\enspace\url{www.mohamedelafrit.com}\enspace\textbullet\enspace CC BY-NC-SA 4.0}}
\fancyfoot[C]{\footnotesize\textcolor{gray}{\thepage}}
\fancyfoot[R]{\footnotesize\textcolor{gray}{2025–2026}}
\renewcommand{\footrulewidth}{0.2pt}

% Style pour la page de garde (pas d'en-tête)
\fancypagestyle{pagegarde}{
    \fancyhf{}
    \renewcommand{\headrulewidth}{0pt}
    \fancyfoot[C]{\footnotesize\textcolor{gray}{\thepage}}
    \renewcommand{\footrulewidth}{0pt}
}

% ============================================================================
% PAGE DE GARDE — \pagegarde{titre}{jour}{theme}
% ============================================================================
\newcommand{\pagegarde}[3]{%
    \thispagestyle{pagegarde}
    \begin{center}
        \vspace*{1.5cm}

        % Logo / Bandeau institution
        {\Large\sffamily\bfseries\textcolor{ipssiBleu}{IPSSI}}\par
        \vspace{2pt}
        {\normalsize\sffamily 2\textsuperscript{e} année Bachelor Informatique}\par

        \vspace{0.8cm}
        {\color{ipssiBleu}\rule{0.6\textwidth}{1.5pt}}\par
        \vspace{0.8cm}

        % Titre du module
        {\LARGE\sffamily\bfseries Mathématiques appliquées\\[4pt] à l'Intelligence Artificielle}\par

        \vspace{1cm}

        % Titre du document
        \begin{tcolorbox}[
            enhanced,
            colback=ipssiBleu!5,
            colframe=ipssiBleu!50,
            boxrule=1pt,
            arc=4pt,
            width=0.85\textwidth,
            halign=center,
            top=10pt, bottom=10pt,
        ]
            {\huge\sffamily\bfseries\textcolor{ipssiBleu!80!black}{#1}}\par
            \vspace{6pt}
            {\Large\sffamily Jour #2 — #3}
        \end{tcolorbox}

        \vspace{1.5cm}

        % Informations auteur
        {\large\sffamily
            \textbf{Auteur :} Mohamed EL AFRIT\\[4pt]
            \url{www.mohamedelafrit.com}\\[8pt]
            \textbf{Année académique :} 2025–2026
        }\par

        \vspace{1.5cm}

        % Licence Creative Commons
        {\color{ipssiBleu}\rule{0.4\textwidth}{0.5pt}}\par
        \vspace{8pt}
        {\small\sffamily
            \textcopyright\ 2025 Mohamed EL AFRIT — IPSSI\\[2pt]
            Ce contenu est distribué sous licence\\[2pt]
            \textbf{Creative Commons BY-NC-SA 4.0}\\[2pt]
            \url{https://creativecommons.org/licenses/by-nc-sa/4.0/deed.fr}
        }\par
    \end{center}
    \newpage
}

% ============================================================================
% NIVEAUX D'EXERCICES
% ============================================================================

% Étoile compatible pdflatex
\newcommand{\starfull}[1]{\textcolor{#1}{$\bigstar$}}
\newcommand{\starempty}{\textcolor{gray!40}{$\bigstar$}}

% Fondamental (1 étoile)
\newcommand{\exofondamental}{%
    \starfull{ipssiVert}\starempty\starempty\enspace%
    {\small\sffamily\textcolor{ipssiVert}{\textbf{Fondamental}}}%
}

% Intermédiaire (2 étoiles)
\newcommand{\exointermediaire}{%
    \starfull{ipssiOrange}\starfull{ipssiOrange}\starempty\enspace%
    {\small\sffamily\textcolor{ipssiOrange}{\textbf{Interm\'{e}diaire}}}%
}

% Avancé (3 étoiles)
\newcommand{\exoavance}{%
    \starfull{ipssiRouge}\starfull{ipssiRouge}\starfull{ipssiRouge}\enspace%
    {\small\sffamily\textcolor{ipssiRouge}{\textbf{Avanc\'{e}}}}%
}

% Commande d'exercice numéroté
\newcounter{exocount}[section]
\newcommand{\exercice}[1]{%
    \stepcounter{exocount}%
    \vspace{8pt}%
    \noindent\textbf{\sffamily Exercice \theexocount}\enspace — \enspace #1\par
    \vspace{4pt}%
}

% ============================================================================
% COMMANDES UTILITAIRES
% ============================================================================

% Séparateur visuel entre sections
\newcommand{\separateur}{%
    \vspace{6pt}
    \begin{center}
        {\color{ipssiBleu!30}\rule{0.3\textwidth}{0.5pt}}
    \end{center}
    \vspace{6pt}
}

% Mot-clé en gras coloré
\newcommand{\motcle}[1]{\textbf{\textcolor{ipssiBleu!80!black}{#1}}}

% Formule encadrée importante
\newcommand{\formule}[1]{%
    \begin{center}
        \fcolorbox{ipssiBleu!50}{ipssiBleu!5}{%
            \quad$\displaystyle #1$\quad%
        }
    \end{center}
}

% ============================================================================
% FIN DU PRÉAMBULE
% ============================================================================


\jourNumero{3}
\jourTheme{Classification}

\begin{document}

\pagegarde{Fiche d'exercices}{3}{Probabilit\'{e}s et classification}

\setstretch{1.08}
\setlength{\parskip}{3pt}
\setlength{\abovedisplayskip}{6pt}
\setlength{\belowdisplayskip}{6pt}

% ======================================================================
% RAPPELS
% ======================================================================
\begin{tcolorbox}[
    enhanced, colback=ipssiCyan!5, colframe=ipssiCyan!50,
    boxrule=0.8pt, arc=4pt,
    title={\sffamily\bfseries Rappels --- Formules du Jour 3},
    coltitle=white, colbacktitle=ipssiCyan!70,
]
\small
\textbf{Probabilit\'{e} conditionnelle :}
$P(A \mid B) = \frac{P(A \cap B)}{P(B)}$
\hfill
\textbf{Perceptron :}
$\yhat = \text{signe}(\vw^T \vx + b)$

\medskip
\textbf{R\`{e}gle de d\'{e}cision :}
$\yhat = 1$ si $g(\vx) = \vw^T\vx + b > 0$, sinon $\yhat = 0$.
\hfill
\textbf{Fronti\`{e}re :} $\vw^T\vx + b = 0$

\medskip
\textbf{Mise \`{a} jour (si erreur) :}
$\vw \leftarrow \vw + \alpha(y_i - \yhat_i)\vx_i$ \;;\;
$b \leftarrow b + \alpha(y_i - \yhat_i)$

\medskip
\textbf{Accuracy :}
$\text{Acc} = \frac{\text{nombre de pr\'{e}dictions correctes}}{n}$
\hfill
\textbf{Dataset :} 30 dev, seuil 45\,k\texteuro, classes \'{e}quilibr\'{e}es 15/15.
\end{tcolorbox}

% ======================================================================
\section*{Exercices crayon-papier --- Niveau fondamental}
% ======================================================================

% ----------------------------------------------------------------------
\exercice{\exofondamental\quad --- Probabilit\'{e}s simples sur un mini-dataset}

Voici un mini-dataset de 8 d\'{e}veloppeurs :

\begin{center}
\renewcommand{\arraystretch}{1.15}
\begin{tabular}{c|c|c|c}
    \toprule
    \textbf{Dev} & \textbf{Exp\'{e}rience} & \textbf{Salaire (k\texteuro)} & \textbf{salaire\_eleve} \\
    \midrule
    A & 2  & 30 & 0 \\
    B & 4  & 42 & 0 \\
    C & 6  & 48 & 1 \\
    D & 8  & 55 & 1 \\
    E & 3  & 35 & 0 \\
    F & 10 & 60 & 1 \\
    G & 1  & 28 & 0 \\
    H & 7  & 46 & 1 \\
    \bottomrule
\end{tabular}
\end{center}

\begin{enumerate}[label=\alph*)]
    \item Calculez $P(\text{salaire\_eleve} = 1)$.
    \item Calculez $P(\text{salaire\_eleve} = 0)$.
        V\'{e}rifiez que $P(0) + P(1) = 1$.
    \item Calculez $P(\text{exp\'{e}rience} \geq 6)$.
    \item Combien de d\'{e}veloppeurs ont \`{a} la fois exp\'{e}rience $\geq 6$ \textbf{et} salaire\_eleve $= 1$ ?
\end{enumerate}

% ----------------------------------------------------------------------
\exercice{\exofondamental\quad --- Probabilit\'{e} conditionnelle \`{a} la main}

En utilisant le m\^{e}me mini-dataset de l'exercice~1 :

\begin{enumerate}[label=\alph*)]
    \item Calculez $P(\text{salaire\_eleve} = 1 \mid \text{exp} \geq 6)$ en utilisant la formule :
    \[
        P(A \mid B) = \frac{\text{nombre de dev avec exp} \geq 6 \textbf{ et } \text{salaire\_eleve} = 1}{\text{nombre de dev avec exp} \geq 6}
    \]

    \item Calculez $P(\text{salaire\_eleve} = 1 \mid \text{exp} < 6)$.

    \item Comparez les deux r\'{e}sultats. Que constatez-vous ?
        Avoir plus de 6 ans d'exp\'{e}rience change-t-il la probabilit\'{e} d'\^{e}tre bien pay\'{e} ?

    \item Calculez le \textbf{ratio} :
    $\frac{P(\text{sal\_eleve}=1 \mid \text{exp}\geq 6)}{P(\text{sal\_eleve}=1 \mid \text{exp}<6)}$.
    Interpr\'{e}tez.
\end{enumerate}

% ----------------------------------------------------------------------
\exercice{\exofondamental\quad --- R\`{e}gle de d\'{e}cision sur 3 observations}

On dispose d'un classifieur lin\'{e}aire avec les poids
$\vw = \begin{pmatrix} 0{,}5 \\ 0{,}8 \end{pmatrix}$ et le biais $b = -4$.

Les features sont : $x_1 = $ exp\'{e}rience, $x_2 = $ nb\_langages.

\begin{enumerate}[label=\alph*)]
    \item Pour les 3 d\'{e}veloppeurs suivants, calculez le score $g(\vx) = \vw^T\vx + b$
        et la pr\'{e}diction $\yhat$ :

    \begin{center}
    \renewcommand{\arraystretch}{1.3}
    \begin{tabular}{c|c|c|c|c|c}
        \toprule
        Dev & $x_1$ & $x_2$ & $g(\vx) = 0{,}5 x_1 + 0{,}8 x_2 - 4$ & $\yhat$ & $y$ (vrai) \\
        \midrule
        1 & 2  & 3 & ? & ? & 0 \\
        2 & 10 & 6 & ? & ? & 1 \\
        3 & 5  & 4 & ? & ? & 0 \\
        \bottomrule
    \end{tabular}
    \end{center}

    \item Quelles pr\'{e}dictions sont correctes ? Calculez l'accuracy.

    \item Le dev~3 a $g(\vx) > 0$ mais $y = 0$. Comment appelle-t-on ce type d'erreur ?
\end{enumerate}

% ----------------------------------------------------------------------
\exercice{\exofondamental\quad --- Tracer une fronti\`{e}re de d\'{e}cision}

On consid\`{e}re le classifieur de l'exercice~3 : $g(\vx) = 0{,}5\,x_1 + 0{,}8\,x_2 - 4$.

\begin{enumerate}[label=\alph*)]
    \item La fronti\`{e}re de d\'{e}cision est l'ensemble des points o\`{u} $g(\vx) = 0$.
        \'{E}crivez l'\'{e}quation de la droite sous la forme $x_2 = a\,x_1 + c$.

    \item Calculez deux points de cette droite :
        quand $x_1 = 0$, que vaut $x_2$ ? Quand $x_1 = 8$ ?

    \item Sur un graphique papier avec $x_1$ en abscisse (0 \`{a} 12) et $x_2$ en ordonn\'{e}e (0 \`{a} 8),
        tracez la droite et les 3 d\'{e}veloppeurs de l'exercice~3.
        Indiquez les zones ``classe 0'' et ``classe 1''.

    \item Le dev~3 est-il du ``bon c\^{o}t\'{e}'' de la fronti\`{e}re ?
\end{enumerate}

% ======================================================================
\section*{Exercices crayon-papier --- Niveau interm\'{e}diaire}
% ======================================================================

% ----------------------------------------------------------------------
\exercice{\exointermediaire\quad --- Trois it\'{e}rations du perceptron \`{a} la main}

On entra\^{i}ne un perceptron avec 2 features, $\alpha = 1$, poids initiaux $\vw = (0, 0)^T$, biais $b = 0$.

Donn\'{e}es d'entra\^{i}nement :

\begin{center}
\renewcommand{\arraystretch}{1.2}
\begin{tabular}{c|cc|c}
    \toprule
    $i$ & $x_1$ & $x_2$ & $y$ \\
    \midrule
    1 & 3 & 4 & 1 \\
    2 & 1 & 1 & 0 \\
    3 & 5 & 2 & 1 \\
    \bottomrule
\end{tabular}
\end{center}

Pour chaque observation (dans l'ordre $i = 1, 2, 3$) :

\begin{enumerate}[label=\alph*)]
    \item Calculez $g(\vx_i) = \vw^T\vx_i + b$ et $\yhat_i = \text{signe}(g)$.

    \item Comparez $\yhat_i$ et $y_i$. Y a-t-il erreur ?

    \item Si oui, calculez $\vw_{\text{new}} = \vw + \alpha(y_i - \yhat_i)\vx_i$
        et $b_{\text{new}} = b + \alpha(y_i - \yhat_i)$.

    \item Remplissez le tableau complet :

    \begin{center}
    \renewcommand{\arraystretch}{1.3}
    \begin{tabular}{c|c|c|c|c|c|c}
        \toprule
        $i$ & $\vw$ avant & $b$ & $g(\vx_i)$ & $\yhat_i$ & Erreur ? & $\vw$ apr\`{e}s \\
        \midrule
        1 & $(0, 0)$ & 0 & ? & ? & ? & ? \\
        2 & ?        & ? & ? & ? & ? & ? \\
        3 & ?        & ? & ? & ? & ? & ? \\
        \bottomrule
    \end{tabular}
    \end{center}
\end{enumerate}

% ----------------------------------------------------------------------
\exercice{\exointermediaire\quad --- Pourquoi le XOR n'est pas lin\'{e}airement s\'{e}parable}

Le probl\`{e}me XOR est d\'{e}fini par :

\begin{center}
\renewcommand{\arraystretch}{1.15}
\begin{tabular}{cc|c}
    \toprule
    $x_1$ & $x_2$ & $y$ \\
    \midrule
    0 & 0 & 0 \\
    0 & 1 & 1 \\
    1 & 0 & 1 \\
    1 & 1 & 0 \\
    \bottomrule
\end{tabular}
\end{center}

\begin{enumerate}[label=\alph*)]
    \item Placez les 4 points sur un graphique ($x_1$ en abscisse, $x_2$ en ordonn\'{e}e).
        Utilisez des couleurs diff\'{e}rentes pour les classes 0 et 1.

    \item Essayez de tracer une \textbf{droite} qui s\'{e}pare les rouges des verts.
        Y arrivez-vous ?

    \item Supposons un classifieur lin\'{e}aire $g(\vx) = w_1 x_1 + w_2 x_2 + b$.
        Montrez par l'absurde qu'il est impossible de trouver $(w_1, w_2, b)$
        qui classifie correctement les 4 points.

        \emph{Indice :} \'{e}crivez les 4 in\'{e}galit\'{e}s que les poids devraient satisfaire :
        $g(0,0) \leq 0$, $g(0,1) > 0$, $g(1,0) > 0$, $g(1,1) \leq 0$.
        Montrez la contradiction.

    \item Pourquoi dit-on que ce probl\`{e}me n'est \textbf{pas lin\'{e}airement s\'{e}parable} ?
        Que faudrait-il pour le r\'{e}soudre ?
\end{enumerate}

% ----------------------------------------------------------------------
\exercice{\exointermediaire\quad --- Calculer l'accuracy d'un classifieur}

Un classifieur a produit les r\'{e}sultats suivants sur 10 observations :

\begin{center}
\renewcommand{\arraystretch}{1.15}
\begin{tabular}{c|cccccccccc}
    \toprule
    $i$       & 1 & 2 & 3 & 4 & 5 & 6 & 7 & 8 & 9 & 10 \\
    \midrule
    $y_i$ (vrai)  & 0 & 1 & 1 & 0 & 1 & 0 & 1 & 0 & 1 & 0 \\
    $\yhat_i$ (pr\'{e}dit)& 0 & 1 & 0 & 0 & 1 & 1 & 1 & 0 & 0 & 0 \\
    \bottomrule
\end{tabular}
\end{center}

\begin{enumerate}[label=\alph*)]
    \item Pour chaque observation, indiquez si la pr\'{e}diction est correcte
        (\checkmark) ou incorrecte (\texttimes).

    \item Calculez l'\textbf{accuracy} : $\text{Acc} = \frac{\text{nb correct}}{n}$.

    \item Combien y a-t-il de \textbf{faux positifs} (pr\'{e}dit 1, vrai 0)
        et de \textbf{faux n\'{e}gatifs} (pr\'{e}dit 0, vrai 1) ?

    \item Si on pr\'{e}disait \textbf{toujours} la classe 0 pour tout le monde,
        quelle serait l'accuracy ? Ce classifieur ``na\"{i}f'' serait-il bon ?
\end{enumerate}

% ======================================================================
\section*{Exercices Python guid\'{e}s --- Niveau interm\'{e}diaire}
% ======================================================================

% ----------------------------------------------------------------------
\exercice{\exointermediaire\quad --- Probabilit\'{e}s empiriques sur le dataset \hfill\normalfont\small\textsf{[Python guid\'{e}]}}

\begin{pythonbox}[Squelette --- Probabilit\'{e}s]
\begin{lstlisting}
import numpy as np

data = np.array([
    [ 0,2,3, 150, 80,28.0],[ 1,1,2,  50,100,26.5],
    [ 1,3,4, 800, 40,35.0],[ 2,2,3, 200, 60,32.0],
    [ 2,4,4,3000, 20,38.0],[ 0,1,1,  30, 50,25.0],
    [ 3,3,3, 100, 80,34.5],[ 1,5,4,  15,100,42.0],
    [ 4,3,3, 500, 60,38.5],[ 5,4,4,1200, 40,44.0],
    [ 5,3,3, 300, 80,41.0],[ 6,5,4,2000, 30,48.0],
    [ 6,4,3, 150, 90,43.5],[ 7,5,4, 800, 50,50.0],
    [ 7,3,2,  80, 70,39.0],[ 4,6,5, 400,100,47.0],
    [ 8,4,4,1500, 40,51.0],[ 5,2,3,4000, 10,40.0],
    [ 9,5,4, 600, 60,52.0],[10,6,4,2500, 30,55.0],
    [10,4,3, 200, 80,48.5],[11,5,5,1000, 50,58.0],
    [12,7,4,3500, 20,60.0],[ 9,3,3,  60, 90,42.0],
    [13,6,4, 800, 70,62.0],[11,4,4,5000, 10,54.0],
    [15,8,4,1200, 50,65.0],[17,6,5,2000, 40,68.0],
    [20,7,4, 500, 60,72.0],[18,5,3, 100, 80,58.0]])

exp = data[:, 0]; sal = data[:, 5]
y = (sal > 45).astype(int)

# 1. P(salaire_eleve = 1)
p1 = ______               # A COMPLETER
print(f"P(sal>45k) = {p1:.2f}")

# 2. P(sal>45k | exp > 10)
mask_exp10 = ______        # A COMPLETER : exp > 10
p_cond = np.mean(______)   # A COMPLETER : y filtre
print(f"P(sal>45k | exp>10) = {p_cond:.2f}")

# 3. P(sal>45k | nb_langages >= 5)
lang = data[:, 1]
mask_lang5 = ______         # A COMPLETER
p_lang = np.mean(______)    # A COMPLETER
print(f"P(sal>45k | lang>=5) = {p_lang:.2f}")

# 4. Quelle feature est la plus "predictive" ?
print("\n--- Comparaison ---")
for seuil_exp in [3, 5, 8, 10]:
    mask = exp > seuil_exp
    if np.sum(mask) > 0:
        p = np.mean(y[mask])
        print(f"P(sal>45k | exp>{seuil_exp:2d}) = {p:.2f} "
              f"({np.sum(mask)} dev)")
\end{lstlisting}
\end{pythonbox}

% ----------------------------------------------------------------------
\exercice{\exointermediaire\quad --- R\`{e}gle de d\'{e}cision et fronti\`{e}re \hfill\normalfont\small\textsf{[Python guid\'{e}]}}

\begin{pythonbox}[Squelette --- Fronti\`{e}re de d\'{e}cision]
\begin{lstlisting}
import numpy as np
import matplotlib.pyplot as plt

# Dataset (experience, nb_langages, salaire)
data = np.array([
    [ 0,2,28.0],[ 1,1,26.5],[ 1,3,35.0],[ 2,2,32.0],
    [ 2,4,38.0],[ 0,1,25.0],[ 3,3,34.5],[ 1,5,42.0],
    [ 4,3,38.5],[ 5,4,44.0],[ 5,3,41.0],[ 6,5,48.0],
    [ 6,4,43.5],[ 7,5,50.0],[ 7,3,39.0],[ 4,6,47.0],
    [ 8,4,51.0],[ 5,2,40.0],[ 9,5,52.0],[10,6,55.0],
    [10,4,48.5],[11,5,58.0],[12,7,60.0],[ 9,3,42.0],
    [13,6,62.0],[11,4,54.0],[15,8,65.0],[17,6,68.0],
    [20,7,72.0],[18,5,58.0]])

X = data[:, :2]  # experience, nb_langages
y = (data[:, 2] > 45).astype(int)

# Poids manuels
w = np.array([______, ______])  # A COMPLETER (essayer 0.5, 0.8)
b = ______                       # A COMPLETER (essayer -4)

# 1. Calculer g(x) pour chaque developpeur
scores = ______                  # A COMPLETER : X @ w + b

# 2. Predictions
y_pred = (scores > 0).astype(int)

# 3. Accuracy
acc = np.mean(______)            # A COMPLETER
print(f"Accuracy = {acc:.0%}")

# 4. Graphique avec frontiere
fig, ax = plt.subplots(figsize=(8, 6))
for c, color, label in [(0,'#e74c3c','<= 45k'),
                          (1,'#2ecc71','> 45k')]:
    mask = y == c
    ax.scatter(X[mask,0], X[mask,1], c=color,
               s=70, label=label, edgecolors='white')

# Frontiere : w1*x1 + w2*x2 + b = 0
#          => x2 = -(w1/w2)*x1 - b/w2
x1_line = np.linspace(-1, 21, 100)
x2_line = ______     # A COMPLETER : equation de la droite
ax.plot(x1_line, x2_line, 'k--', lw=2, label='Frontiere')
ax.set_xlim(-1, 21); ax.set_ylim(0, 9)
ax.set_xlabel("Experience"); ax.set_ylabel("Nb langages")
ax.set_title(f"Frontiere (accuracy={acc:.0%})")
ax.legend(); ax.grid(True, alpha=0.3)
plt.tight_layout(); plt.show()
\end{lstlisting}
\end{pythonbox}

% ----------------------------------------------------------------------
\exercice{\exointermediaire\quad --- Perceptron from scratch \hfill\normalfont\small\textsf{[Python guid\'{e}]}}

\begin{pythonbox}[Squelette --- Perceptron]
\begin{lstlisting}
import numpy as np
import matplotlib.pyplot as plt

# Dataset (meme que exercice 9)
data = np.array([
    [ 0,2,28.0],[ 1,1,26.5],[ 1,3,35.0],[ 2,2,32.0],
    [ 2,4,38.0],[ 0,1,25.0],[ 3,3,34.5],[ 1,5,42.0],
    [ 4,3,38.5],[ 5,4,44.0],[ 5,3,41.0],[ 6,5,48.0],
    [ 6,4,43.5],[ 7,5,50.0],[ 7,3,39.0],[ 4,6,47.0],
    [ 8,4,51.0],[ 5,2,40.0],[ 9,5,52.0],[10,6,55.0],
    [10,4,48.5],[11,5,58.0],[12,7,60.0],[ 9,3,42.0],
    [13,6,62.0],[11,4,54.0],[15,8,65.0],[17,6,68.0],
    [20,7,72.0],[18,5,58.0]])

X = data[:, :2]
y = (data[:, 2] > 45).astype(int)
n = len(y)

# Initialisation
w = np.zeros(2)
b = 0.0
alpha = ______       # A COMPLETER (essayer 0.1)
n_epochs = 100

for epoch in range(n_epochs):
    n_errors = 0
    for i in range(n):
        # 1. Score
        z = ______   # A COMPLETER : w @ X[i] + b

        # 2. Prediction
        y_pred = ______  # A COMPLETER : 1 si z>0

        # 3. Mise a jour si erreur
        if y_pred != y[i]:
            error = ______  # A COMPLETER : y[i] - y_pred
            w = w + ______  # A COMPLETER
            b = b + ______  # A COMPLETER
            n_errors += 1

    if epoch < 5 or epoch % 20 == 19:
        acc = np.mean([(1 if (w@X[j]+b)>0 else 0)==y[j]
                        for j in range(n)])
        print(f"Epoch {epoch+1:3d} : "
              f"erreurs={n_errors}, acc={acc:.0%}")

# Graphique final
fig, ax = plt.subplots(figsize=(8, 6))
for c, color, label in [(0,'#e74c3c','<= 45k'),
                          (1,'#2ecc71','> 45k')]:
    mask = y == c
    ax.scatter(X[mask,0], X[mask,1], c=color,
               s=70, label=label, edgecolors='white')

if abs(w[1]) > 1e-6:
    x1 = np.linspace(-1, 21, 100)
    x2 = ______  # A COMPLETER : frontiere
    ax.plot(x1, x2, 'k--', lw=2, label='Frontiere')

ax.set_xlim(-1, 21); ax.set_ylim(0, 9)
ax.set_xlabel("Experience"); ax.set_ylabel("Nb langages")
ax.set_title("Perceptron from scratch")
ax.legend(); ax.grid(True, alpha=0.3)
plt.tight_layout(); plt.show()
\end{lstlisting}
\end{pythonbox}

% ======================================================================
\section*{Exercices Python autonomes --- Niveau avanc\'{e}}
% ======================================================================

% ----------------------------------------------------------------------
\exercice{\exoavance\quad --- Perceptron from scratch vs scikit-learn \hfill\normalfont\small\textsf{[Python autonome]}}

\'{E}crivez un programme Python complet qui :

\begin{enumerate}[label=\alph*)]
    \item Charge le dataset (2 features : exp\'{e}rience, nb\_langages) et la cible binaire (seuil 45\,k\texteuro).

    \item \textbf{M\'{e}thode 1 --- From scratch :} impl\'{e}mentez le perceptron
        ($\alpha = 0{,}1$, 100 \'{e}poques). Enregistrez l'accuracy \`{a} chaque \'{e}poque.

    \item \textbf{M\'{e}thode 2 --- scikit-learn :}
        \texttt{from sklearn.linear\_model import Perceptron}.
        Entra\^{i}nez avec les m\^{e}mes donn\'{e}es.

    \item Affichez un \textbf{tableau comparatif} :
    \begin{center}
    \renewcommand{\arraystretch}{1.2}
    \begin{tabular}{l|c|c|c|c}
        \toprule
        \textbf{M\'{e}thode} & $w_1$ & $w_2$ & $b$ & Accuracy \\
        \midrule
        From scratch & ? & ? & ? & ? \\
        scikit-learn & ? & ? & ? & ? \\
        \bottomrule
    \end{tabular}
    \end{center}

    \item Tracez un \textbf{unique graphique} montrant :
        le nuage de points color\'{e}, la fronti\`{e}re from scratch (rouge),
        la fronti\`{e}re scikit-learn (verte), une l\'{e}gende avec les accuracies.

    \item Tracez la \textbf{courbe d'accuracy par \'{e}poque} du perceptron from scratch.
        Combien d'\'{e}poques sont n\'{e}cessaires pour converger ?
\end{enumerate}

\begin{attentionbox}[Consignes]
    Utilisez NumPy, Matplotlib et scikit-learn. Commentez en fran\c{c}ais.
    Les deux fronti\`{e}res ne seront probablement pas identiques : expliquez pourquoi.
\end{attentionbox}

% ----------------------------------------------------------------------
\exercice{\exoavance\quad --- Perceptron sur XOR + solution \hfill\normalfont\small\textsf{[Python autonome]}}

\'{E}crivez un programme qui :

\begin{enumerate}[label=\alph*)]
    \item D\'{e}finit le dataset XOR (4 points, 2 classes).

    \item Entra\^{i}ne un perceptron from scratch (200 \'{e}poques, $\alpha = 0{,}1$).
        Enregistrez le nombre d'erreurs par \'{e}poque.

    \item Tracez \textbf{deux graphiques} c\^{o}te \`{a} c\^{o}te :
    \begin{itemize}
        \item \`{A} gauche : les 4 points du XOR color\'{e}s + la derni\`{e}re fronti\`{e}re
            tent\'{e}e par le perceptron.
        \item \`{A} droite : le nombre d'erreurs par \'{e}poque (montrer l'oscillation).
    \end{itemize}

    \item Concluez : le perceptron converge-t-il ? Pourquoi ?

    \item \textbf{Bonus --- Solution avec transformation de features :}
        Ajoutez une 3\textsuperscript{e} feature $x_3 = x_1 \times x_2$ (le produit).
        Le nouveau dataset est :

    \begin{center}
    \renewcommand{\arraystretch}{1.1}
    \begin{tabular}{ccc|c}
        \toprule
        $x_1$ & $x_2$ & $x_3 = x_1 x_2$ & $y$ \\
        \midrule
        0 & 0 & 0 & 0 \\
        0 & 1 & 0 & 1 \\
        1 & 0 & 0 & 1 \\
        1 & 1 & 1 & 0 \\
        \bottomrule
    \end{tabular}
    \end{center}

        Entra\^{i}nez le perceptron sur ce dataset augment\'{e}.
        Converge-t-il cette fois ? Tracez le r\'{e}sultat.
\end{enumerate}

\begin{intuitionbox}[Indice]
    La transformation $x_3 = x_1 x_2$ rend le probl\`{e}me \textbf{lin\'{e}airement s\'{e}parable}
    dans $\mathbb{R}^3$. C'est l'id\'{e}e fondamentale des \textbf{r\'{e}seaux de neurones} :
    les couches cach\'{e}es cr\'{e}ent automatiquement de nouvelles features.
\end{intuitionbox}

% ======================================================================
\vfill
\separateur

\begin{center}
    \small\sffamily
    \textbf{R\'{e}partition :}\quad
    4 exercices \starfull{ipssiVert} fondamentaux \quad
    3 exercices \starfull{ipssiOrange}\starfull{ipssiOrange} interm\'{e}diaires \quad
    3 guid\'{e}s Python \starfull{ipssiOrange}\starfull{ipssiOrange} \quad
    2 autonomes \starfull{ipssiRouge}\starfull{ipssiRouge}\starfull{ipssiRouge}
\end{center}

\begin{center}
    {\footnotesize\sffamily\textcolor{gray}{%
        \textcopyright\ 2025 Mohamed EL AFRIT --- IPSSI \enspace\textbullet\enspace
        \url{www.mohamedelafrit.com} \enspace\textbullet\enspace
        CC BY-NC-SA 4.0
    }}
\end{center}

\end{document}
